\chapter{Conclusion}\label{chap:conclusion}

L'atelier~3 a transformé une application MVC2 en application JSF sans toucher à une seule classe de la couche métier ni de la
persistance. Cette observation est la plus importante du travail : la séparation en couches de l'atelier~2 a rempli
son rôle, et le coût d'un changement de technologie web s'est limité à quinze classes et seize pages.

Trois idées en ressortent. D'abord, JSF déplace la réflexion du \textit{chemin d'une requête} vers le \textit{cycle de vie d'une vue} ; les
portées (\code{@ViewScoped}, jamais \code{@SessionScoped} ici) et le motif Post/Redirect/Get demandent une attention
constante. Ensuite, une politique de sécurité stricte contraint le choix des composants : renoncer à
\code{h:commandLink} et à \code{f:ajax} est un arbitrage, pas une limite. Enfin, sortir la logique d'affichage dans des
méthodes Java rend les vues plus simples, les tests possibles et les analyseurs statiques silencieux.

Les limites sont connues : la commande n'est pas traitée, le limiteur de connexions est local à un serveur, et
Tomcat, MySQL et Maven n'ont été vérifiés que sur ma machine, les tests automatisés couvrant ce qui s'exécute sans serveur. Le rapport
rend compte de ces deux niveaux de vérification et des incidents qui les ont séparés.
